\documentclass[12pt,a4paper]{article}

% === Kodowanie i język ===
\usepackage[utf8]{inputenc}
\usepackage[T1]{fontenc}
\usepackage[polish]{babel}

% === Układ strony ===
\usepackage[margin=2.5cm, headheight=14.5pt]{geometry}
\usepackage{setspace}
\onehalfspacing

% === Grafika i kolory ===
\usepackage{graphicx}
\usepackage{xcolor}
\usepackage{tikz}
\usetikzlibrary{arrows.meta, positioning, shapes.geometric, calc}

% === Tabele ===
\usepackage{booktabs}
\usepackage{longtable}
\usepackage{array}
\usepackage{multirow}

% === Listingi kodu ===
\usepackage{listings}
\definecolor{codebg}{RGB}{248,248,248}
\definecolor{codegreen}{RGB}{40,160,40}
\definecolor{codegray}{RGB}{128,128,128}
\definecolor{codepurple}{RGB}{160,32,240}
\definecolor{codeblue}{RGB}{0,80,180}

\lstdefinestyle{python}{
    language=Python,
    backgroundcolor=\color{codebg},
    basicstyle=\ttfamily\small,
    keywordstyle=\color{codeblue}\bfseries,
    stringstyle=\color{codegreen},
    commentstyle=\color{codegray}\itshape,
    numberstyle=\tiny\color{codegray},
    numbers=left,
    numbersep=8pt,
    frame=single,
    framerule=0.4pt,
    rulecolor=\color{codegray!50},
    breaklines=true,
    showstringspaces=false,
    tabsize=4,
    captionpos=b,
    aboveskip=10pt,
    belowskip=10pt,
    morekeywords={self, True, False, None, as, assert, with},
    literate=
        {ą}{{\k{a}}}1 {ć}{{\'c}}1 {ę}{{\k{e}}}1 {ł}{{\l{}}}1
        {ń}{{\'n}}1 {ó}{{\'o}}1 {ś}{{\'s}}1 {ź}{{\'z}}1 {ż}{{\.z}}1
        {Ą}{{\k{A}}}1 {Ć}{{\'C}}1 {Ę}{{\k{E}}}1 {Ł}{{\L{}}}1
        {Ń}{{\'N}}1 {Ó}{{\'O}}1 {Ś}{{\'S}}1 {Ź}{{\'Z}}1 {Ż}{{\.Z}}1,
}
\lstset{style=python}

\lstdefinestyle{yaml}{
    basicstyle=\ttfamily\small,
    backgroundcolor=\color{codebg},
    frame=single,
    framerule=0.4pt,
    rulecolor=\color{codegray!50},
    breaklines=true,
    showstringspaces=false,
    captionpos=b,
    aboveskip=10pt,
    belowskip=10pt,
}

\lstdefinestyle{tree}{
    basicstyle=\ttfamily\small,
    backgroundcolor=\color{codebg},
    frame=single,
    framerule=0.4pt,
    rulecolor=\color{codegray!50},
    breaklines=true,
    captionpos=b,
    aboveskip=10pt,
    belowskip=10pt,
}

% === Matematyka ===
\usepackage{amsmath}
\usepackage{amssymb}
\usepackage{braket}

% === Cudzysłowy ===
\usepackage{csquotes}

% === Hiperłącza ===
\usepackage[colorlinks=true, linkcolor=blue!70!black, citecolor=blue!70!black, urlcolor=blue!60!black]{hyperref}

% === Bibliografia ===
\usepackage[backend=biber, style=numeric, sorting=none, urldate=long]{biblatex}
\addbibresource{raport.bib}

% === Nagłówek i stopka ===
\usepackage{fancyhdr}
\pagestyle{fancy}
\fancyhf{}
\fancyhead[L]{\small\textit{IDUB\_2026 --- Raport o stanie projektu}}
\fancyhead[R]{\small\thepage}
\renewcommand{\headrulewidth}{0.4pt}

% === Formatowanie sekcji ===
\usepackage{titlesec}
\titleformat{\section}{\Large\bfseries}{\thesection.}{0.5em}{}
\titleformat{\subsection}{\large\bfseries}{\thesubsection.}{0.5em}{}
\titleformat{\subsubsection}{\normalsize\bfseries}{\thesubsubsection.}{0.5em}{}

% === Otoczenia informacyjne ===
\usepackage{tcolorbox}
\tcbuselibrary{breakable, skins}

\newtcolorbox{infobox}[1][]{
    colback=blue!5, colframe=blue!40!black,
    fonttitle=\bfseries, title=#1,
    breakable, enhanced
}

\newtcolorbox{warningbox}[1][]{
    colback=red!5, colframe=red!60!black,
    fonttitle=\bfseries, title=#1,
    breakable, enhanced
}

\newtcolorbox{tipbox}[1][]{
    colback=green!5, colframe=green!40!black,
    fonttitle=\bfseries, title=#1,
    breakable, enhanced
}

% === Metadane ===
\title{
    \vspace{-1cm}
    {\Large Raport o stanie projektu}\\[0.5em]
    {\LARGE\bfseries IDUB\_2026}\\[0.3em]
    {\large Głębokie uczenie do klasyfikacji i rekonstrukcji\\kwantowych stanów optycznych}
}
\author{
    Audyt kodu i plan działań\\[0.3em]
    \small Repozytorium: \url{https://github.com/e-mcQuantum-AI/IDUB_2026}
}
\date{28 marca 2026}

% =====================================================================
\begin{document}
\maketitle
\thispagestyle{fancy}

\tableofcontents
\newpage

% =====================================================================
\section{Opis projektu}
\label{sec:opis}

Projekt ma na celu stworzenie narzędzia do generowania syntetycznych danych obrazowych (funkcji Wignera) kwantowych stanów optycznych, które posłużą jako zbiór treningowy do klasyfikacji i rekonstrukcji stanów za pomocą głębokiego uczenia.

Obecny kod obejmuje:
\begin{itemize}
    \item Implementację 7 stanów kwantowych (Fock, koherentny, kot Schrödingera, GKP, dwumianowy, termiczny, próżnia),
    \item Kanały szumu kwantowego (straty fotonowe, mieszanina stanów),
    \item Pomiar funkcją Wignera z potokiem pomiarowym (pipeline),
    \item Skrypty do generowania i dzielenia obrazów zbioru danych.
\end{itemize}

\noindent\textbf{Rozmiar kodu}: ok.~330 linii Pythona w~18 plikach.

% =====================================================================
\section{Wprowadzenie do pojęć --- dla początkujących}
\label{sec:pojecia}

Ten rozdział wyjaśnia kluczowe pojęcia z fizyki kwantowej i programowania, które pojawiają się w dalszej części raportu.

% ----- Fizyka kwantowa -----
\subsection{Pojęcia z fizyki kwantowej}

\subsubsection{Stan kwantowy i notacja Diraca}

\textbf{Stan kwantowy} to matematyczny opis układu kwantowego. W~tym projekcie opisujemy stany światła (fotony w~jednym modzie pola elektromagnetycznego).

\textbf{Ket} $\ket{\psi}$ (czytaj: \emph{ket psi}) --- wektor opisujący stan kwantowy w~notacji Diraca. Na przykład $\ket{3}$ oznacza stan z~dokładnie 3 fotonami. W~kodzie ket to wektor kolumnowy (obiekt \texttt{Qobj} z~biblioteki QuTiP).

\textbf{Bra} $\bra{\psi}$ --- wektor sprzężony do ketu (transponowany i~zespolenie sprzężony, czyli hermitowsko sprzężony). Potrzebny do obliczeń, np. iloczynu skalarnego $\braket{\psi|\psi}$.

\textbf{Iloczyn skalarny} $\braket{\psi|\psi}$ --- miara ,,nakładania się'' dwóch stanów. Jeśli $\braket{\psi|\psi} = 1$, stan jest \textbf{znormalizowany} (ma sens fizyczny).

\subsubsection{Stan czysty a stan mieszany}

\begin{itemize}
    \item \textbf{Stan czysty} --- układ jest w~jednym, określonym stanie kwantowym $\ket{\psi}$. Przykład: dokładnie 3 fotony.
    \item \textbf{Stan mieszany} --- układ jest w~jednym z~kilku stanów z~pewnym prawdopodobieństwem (mieszanina statystyczna). Przykład: stan termiczny --- nie wiemy, ile jest fotonów, znamy tylko średnią. Stanu mieszanego \textbf{nie da się} opisać pojedynczym ketem.
\end{itemize}

\subsubsection{Macierz gęstości}

\textbf{Macierz gęstości} $\rho$ --- uniwersalny sposób opisu stanu kwantowego, działa zarówno dla stanów czystych, jak i~mieszanych:
\begin{align}
    \text{Stan czysty:} \quad & \rho = \ket{\psi}\bra{\psi} \label{eq:rho-pure} \\
    \text{Stan mieszany:} \quad & \rho = \sum_i p_i \ket{\psi_i}\bra{\psi_i} \label{eq:rho-mixed}
\end{align}

\noindent W~kodzie: \texttt{state.density\_matrix()} zwraca macierz kwadratową.

\begin{infobox}[Fizyczność macierzy gęstości]
Poprawna macierz gęstości musi spełniać 3 warunki:
\begin{enumerate}
    \item \textbf{Hermitowskość}: $\rho = \rho^\dagger$ --- macierz jest równa swojemu sprzężeniu hermitowskiemu (tzn.\ transpozycji ze sprzężeniem zespolonym).
    \item \textbf{Ślad = 1}: $\mathrm{Tr}(\rho) = 1$ --- prawdopodobieństwa sumują się do~1.
    \item \textbf{Dodatnia półokreśloność}: wszystkie wartości własne $\geq 0$ --- prawdopodobieństwa nie mogą być ujemne.
\end{enumerate}
\end{infobox}

\subsubsection{Przestrzeń Hilberta i cutoff}

\textbf{Przestrzeń Hilberta} to przestrzeń matematyczna, w~której ,,żyją'' stany kwantowe. Jej wymiar jest teoretycznie nieskończony (oscylator może zawierać dowolnie wiele fotonów), ale w~komputerze musimy ją \textbf{obciąć} --- parametr \texttt{cutoff} określa, ile stanów bazowych ($0, 1, 2, \ldots, \texttt{cutoff}-1$ fotonów) uwzględniamy. Na przykład \texttt{cutoff=32} oznacza, że rozważamy stany od~0 do~31 fotonów.

\subsubsection{Operatory kreacji i anihilacji}

Operatory matematyczne na przestrzeni Hilberta:
\begin{align}
    \hat{a}\ket{n} &= \sqrt{n}\,\ket{n-1} && \text{anihilacja --- ,,zabiera'' foton} \\
    \hat{a}^\dagger\ket{n} &= \sqrt{n+1}\,\ket{n+1} && \text{kreacja --- ,,dodaje'' foton} \\
    \hat{n} = \hat{a}^\dagger\hat{a} && && \text{operator liczby fotonów: } \hat{n}\ket{n} = n\ket{n}
\end{align}

\noindent W~kodzie: \texttt{destroy()} = $\hat{a}$, \texttt{create()} = $\hat{a}^\dagger$.

\subsubsection{Funkcja Wignera}

Sposób wizualizacji stanu kwantowego jako ,,mapy cieplnej'' w~\textbf{przestrzeni fazowej} (oś~$x$ = pozycja~$q$, oś~$y$ = pęd~$p$). Formalnie, dla macierzy gęstości~$\rho$ (w~jednostkach, w~których $\hbar = 2$, zgodnie z~konwencją QuTiP):
\begin{equation}
    W(q, p) = \frac{1}{\pi\hbar} \int_{-\infty}^{\infty} \bra{q + y}\rho\ket{q - y}\, e^{2ipy/\hbar}\, dy
    \label{eq:wigner}
\end{equation}
gdzie $\hbar$ to zredukowana stała Plancka (podstawowa stała fizyczna łącząca skalę kwantową z~klasyczną). W~praktyce nie trzeba liczyć tej całki ręcznie --- biblioteka QuTiP oblicza ją automatycznie funkcją \texttt{wigner(rho, xvec, xvec)}.

Każdy typ stanu ma charakterystyczny wzór na wykresie Wignera --- to właśnie te wzory projekt chce klasyfikować za pomocą uczenia maszynowego. Funkcja Wignera może przyjmować \textbf{wartości ujemne} --- to cecha typowo kwantowa, niemożliwa w~fizyce klasycznej (ujemność Wignera jest świadectwem nieklasyczności stanu). Schemat potoku pomiarowego przedstawia Rysunek~\ref{fig:pipeline}.

\subsubsection{Operator przesunięcia}

$\hat{D}(\alpha)$ --- przesuwa stan w~przestrzeni fazowej o~wektor $\alpha$ (liczba zespolona: $\mathrm{Re}(\alpha)$ = przesunięcie w~$q$, $\mathrm{Im}(\alpha)$ = przesunięcie w~$p$). W~kodzie: \texttt{displace(cutoff, alpha)}.

\subsubsection{Równanie Lindblada i operatory kolapsu}

\textbf{Równanie Lindblada} (zwane też \emph{master equation})~\cite{lindblad1976} --- równanie opisujące ewolucję czasową stanu kwantowego oddziałującego z~otoczeniem (układ otwarty). W~tym projekcie służy do symulacji szumu: zamiast izolowanej ewolucji unitarnej, stan $\rho$ ewoluuje z~uwzględnieniem strat do otoczenia.

\textbf{Operator kolapsu} $\hat{c}$ --- operator opisujący konkretny mechanizm oddziaływania z~otoczeniem w~równaniu Lindblada. Różne operatory kolapsu modelują różne rodzaje szumu:
\begin{itemize}
    \item $\hat{c} = \sqrt{\gamma}\,\hat{a}$ --- straty fotonowe (fotony uciekają z~układu),
    \item $\hat{c} = \sqrt{\gamma}\,\hat{a}^\dagger\hat{a}$ --- defazowanie (utrata informacji o~fazie),
    \item $\hat{c} = \sqrt{\gamma}\,\hat{a}^\dagger$ --- spontaniczne wzmocnienie (dodawanie fotonów).
\end{itemize}

\noindent W~kodzie: parametr \texttt{c\_ops} w~funkcji \texttt{mesolve()} z~biblioteki QuTiP~\cite{qutip}.

% ----- Stany kwantowe -----
\subsection{Stany kwantowe zaimplementowane w projekcie}
\label{sec:stany-opis}

\begin{table}[htbp]
\centering
\caption{Stany kwantowe zaimplementowane w projekcie.}
\label{tab:stany}
\begin{tabular}{@{}llp{5.5cm}p{4cm}@{}}
\toprule
\textbf{Stan} & \textbf{Symbol} & \textbf{Opis intuicyjny} & \textbf{Przykład} \\
\midrule
Próżnia & $\ket{0}$ & Brak fotonów, lecz niezerowe fluktuacje kwantowe (energia punktu zerowego). & Najniższy stan energetyczny pola \\
\addlinespace
Fock & $\ket{n}$ & Dokładnie $n$ fotonów. Stan o~określonej liczbie cząstek. & Foton z~emitera kwantowego \\
\addlinespace
Koherentny & $\ket{\alpha}$ & Stan najbliższy klasycznemu światłu. $\langle\hat{n}\rangle = |\alpha|^2$. & Światło lasera \\
\addlinespace
Kot & $\ket{\alpha}+\ket{-\alpha}$ & Superpozycja dwóch stanów koherentnych --- ,,jednocześnie tu i~tam''. & Eksperyment Schrödingera \\
\addlinespace
Termiczny & $\rho_\mathrm{th}$ & Stan mieszany w~równowadze termicznej (rozkład Bosego--Einsteina). & Promieniowanie żarówki \\
\addlinespace
Dwumianowy & $\ket{\psi_\mathrm{bin}}$ & Superpozycja stanów Focka ze współczynnikami dwumianowymi~\cite{michael2016}. & Pamięci kwantowe \\
\addlinespace
GKP & $\ket{\mathrm{GKP}}$ & Kod korekcji błędów: kubit zakodowany w~oscylatorze harmonicznym~\cite{gottesman2001}. & Komputery kwantowe ze zmiennymi ciągłymi \\
\bottomrule
\end{tabular}
\end{table}

% ----- Kanały szumu -----
\subsection{Kanały szumu kwantowego}

Kanał szumu to operacja symulująca niedoskonałości fizyczne --- wprowadza szum do stanu kwantowego. Zaimplementowane w~projekcie kanały (straty, mieszanina) zmieniają stan w~sposób nieodwracalny:

\begin{table}[htbp]
\centering
\caption{Kanały szumu zaimplementowane w projekcie.}
\label{tab:kanaly}
\begin{tabular}{@{}lp{6cm}p{5cm}@{}}
\toprule
\textbf{Kanał} & \textbf{Co robi} & \textbf{Analogia} \\
\midrule
Straty fotonowe & Fotony ,,uciekają'' z~układu z~intensywnością~$\gamma$ (im większe~$\gamma$, tym więcej strat). Zmniejsza $\langle\hat{n}\rangle$. & Światło przez brudne szkło \\
\addlinespace
Mieszanina & Miesza dwa stany: $\rho' = p\rho + (1-p)\rho_\mathrm{other}$. & Niedoskonały detektor \\
\bottomrule
\end{tabular}
\end{table}

% ----- Pojęcia programistyczne -----
\subsection{Pojęcia programistyczne}

\begin{description}
    \item[\texttt{\_\_init\_\_.py}] Specjalny plik w~Pythonie, który mówi interpreterowi: ,,ten katalog jest pakietem Pythona''. Bez niego nie można pisać \texttt{from src.physics import FockState}.
    \item[Klasa abstrakcyjna (ABC)] Klasa definiująca ,,umowę'' (interfejs): mówi, jakie metody muszą mieć klasy potomne, ale sama ich nie implementuje. Np.~\texttt{QuantumState} wymaga metody \texttt{ket()}.
    \item[Dziedziczenie] Mechanizm, w~którym klasa potomna przejmuje cechy klasy nadrzędnej. \texttt{class LossChannel(QuantumChannel)} oznacza: ,,LossChannel jest rodzajem QuantumChannel''.
    \item[Type hints] Adnotacje typów: \texttt{def ket(self) -> Qobj:} mówi ,,ta metoda zwraca Qobj''. Python ich nie wymusza, ale pomagają w~wykrywaniu błędów.
    \item[Docstring] Komentarz dokumentacyjny w~potrójnych cudzysłowach \texttt{"""}. Opisuje co robi funkcja/klasa.
    \item[\texttt{requirements.txt}] Plik z~listą bibliotek (np.~\texttt{qutip>=5.0}). Instalacja: \texttt{pip install -r requirements.txt}.
    \item[pytest / fixture] Framework do testów automatycznych. Fixture to przygotowane dane współdzielone przez testy.
    \item[Pipeline] Wzorzec ,,potok'': dane przechodzą kolejno przez etapy. Tu: stan $\to$ szum $\to$ pomiar Wignera $\to$ obraz.
    \item[CI/CD] Automatyczne uruchamianie testów przy każdym \texttt{git push}.
\end{description}

% =====================================================================
\section{Ocena obecnego stanu}
\label{sec:ocena}

\subsection{Architektura --- ocena: DOBRA}

Projekt ma przejrzystą strukturę obiektową z~użyciem klas abstrakcyjnych (ABC):

\begin{lstlisting}[style=tree, caption={Struktura katalogów projektu.}, label={lst:tree}]
src/physics/
|-- state/          -- implementacje stanow kwantowych
|-- noise/          -- kanaly szumowe
|-- measurement/    -- pomiar funkcja Wignera + pipeline
|-- divider.py      -- narzedzie do dzielenia obrazow
|-- generate.py     -- skrypt generujacy zbior danych
|-- hilbert.py      -- operatory przestrzeni Hilberta
`-- validation.py   -- walidacja stanow fizycznych
\end{lstlisting}

Dobór bibliotek jest trafny: \textbf{QuTiP}~\cite{qutip}, NumPy, SciPy, Matplotlib, Pillow (biblioteka do przetwarzania obrazów, importowana jako~PIL).

Zależności klas przedstawia Rysunek~\ref{fig:architektura}.

\subsection{Poprawność fizyczna --- ocena: DOBRA z zastrzeżeniami}

\begin{table}[htbp]
\centering
\caption{Poprawność implementacji stanów kwantowych.}
\label{tab:poprawnosc}
\begin{tabular}{@{}llc@{}}
\toprule
\textbf{Stan} & \textbf{Implementacja} & \textbf{Poprawność} \\
\midrule
Próżnia (\texttt{vacuum.py})    & \texttt{basis(N, 0)}                 & \checkmark \\
Fock (\texttt{fock.py})          & \texttt{basis(N, n)}                 & \checkmark \\
Koherentny (\texttt{coherent.py}) & QuTiP \texttt{coherent()}          & \checkmark \\
Kot (\texttt{cat.py})            & superpozycja stanów koherentnych     & \checkmark \\
Termiczny (\texttt{thermal.py})  & \texttt{thermal\_dm()}, stan mieszany & \checkmark \\
GKP (\texttt{gkp.py})           & obwiednia gaussowska + przesunięcia  & $\triangle$ \\
Dwumianowy (\texttt{binomal.py}) & współczynniki dwumianowe             & $\triangle$ \\
\bottomrule
\end{tabular}

\smallskip
\footnotesize \checkmark\ = poprawna, \quad $\triangle$ = wymaga weryfikacji
\end{table}

\subsection{Jakość kodu --- ocena: PRZECIĘTNA}

\begin{table}[htbp]
\centering
\caption{Ocena aspektów jakości kodu.}
\label{tab:jakosc}
\begin{tabular}{@{}llp{6cm}@{}}
\toprule
\textbf{Aspekt} & \textbf{Ocena} & \textbf{Uwagi} \\
\midrule
Struktura katalogów & Dobra    & Logiczny podział na moduły \\
Czytelność          & Dobra    & Krótkie, jasne funkcje \\
Styl (PEP8)         & Dobra    & Generalnie zgodny \\
Dokumentacja        & Słaba    & Brak docstringów \\
Typowanie           & Brak     & Zero adnotacji typów \\
Obsługa błędów      & Brak     & Brak walidacji wejścia \\
Testy               & Brak     & Zero testów \\
Konfiguracja        & Słaba    & Ścieżki na sztywno \\
\bottomrule
\end{tabular}
\end{table}

% =====================================================================
\section{Wykryte błędy i problemy}
\label{sec:bledy}

\subsection{Błędy krytyczne}

\begin{warningbox}[B1. Brak plików \texttt{\_\_init\_\_.py}]
Żaden pakiet nie ma pliku \texttt{\_\_init\_\_.py}. Import modułów przez \texttt{from src.physics import ...} jest niemożliwy. Dotyczy: \texttt{src/}, \texttt{src/physics/}, \texttt{src/physics/state/}, \texttt{src/physics/noise/}, \texttt{src/physics/measurement/}.
\end{warningbox}

\begin{warningbox}[B2. Błąd matematyczny w \texttt{generate.py} (linia 44)]
Kod tworzy superpozycję dwóch stanów Focka zamiast prawdziwego stanu dwumianowego:
\begin{lstlisting}[numbers=none]
state_binomial = (fock(N, n) + p * fock(N, n+1)).unit()
\end{lstlisting}
To NIE jest stan dwumianowy --- to $\ket{n} + p\ket{n+1}$. Prawdziwa implementacja istnieje w~\texttt{state/binomal.py}, ale nie jest używana przez generator.
\end{warningbox}

\begin{warningbox}[B3. Literówka w nazwie pliku]
Plik \texttt{state/binomal.py} powinien nazywać się \texttt{binomial.py}.
\end{warningbox}

\begin{warningbox}[B4. Kanały szumu nie dziedziczą po klasie bazowej]
\texttt{LossChannel} i~\texttt{MixtureChannel} nie dziedziczą po \texttt{QuantumChannel} (zdefiniowanej w~\texttt{noise/base.py}). Oznacza to, że Python nie wie, iż te klasy są kanałami kwantowymi --- nie można ich używać zamiennie w~kodzie oczekującym obiektu typu \texttt{QuantumChannel} (np.~w~\texttt{MeasurementPipeline}).
\end{warningbox}

\subsection{Problemy istotne}

\textbf{P1. Ścieżki wpisane na sztywno} --- \texttt{divider.py}: \texttt{'quantum\_dataset/clean/'}, \texttt{generate.py}: \texttt{'quantum\_dataset/'}. Brak przenośności.

\textbf{P2. Brak \texttt{requirements.txt}} --- Zależności (QuTiP, NumPy, SciPy, Matplotlib, Pillow, OpenCV) nie są udokumentowane.

\textbf{P3. Brak walidacji danych wejściowych} --- \texttt{BinomialState}: brak sprawdzenia $p \in [0,1]$, $N \geq 0$; \texttt{MixtureChannel}: brak sprawdzenia $p \in [0,1]$; \texttt{LossChannel}: brak sprawdzenia $\gamma \geq 0$.

\textbf{P4. Wzór GKP nieudokumentowany} --- Stan GKP (Gottesman--Kitaev--Preskill)~\cite{gottesman2001} to kwantowy kod korekcji błędów kodujący kubit w~oscylatorze harmonicznym. Idealny stan GKP to nieskończony grzebień delt Diraca w~przestrzeni pozycji, rozmieszczonych co $2\sqrt{\pi}$. W~praktyce stosuje się przybliżenie z~obwiednią gaussowską o~szerokości~$\delta$. Brak referencji w~kodzie utrudnia weryfikację:
\begin{lstlisting}[numbers=none, caption={Nieudokumentowany wzór GKP w \texttt{gkp.py:25}.}]
weight = np.exp(-(2 * s * np.sqrt(np.pi))**2 / (2 * self.delta**2))
\end{lstlisting}

\textbf{P5. Mieszanie języków} --- Komentarze w~\texttt{generate.py} po polsku, reszta kodu po angielsku.

\subsection{Problemy mniejsze}

\begin{itemize}
    \item \textbf{P6.} Stała magiczna \texttt{-1e-10} w~\texttt{validation.py} --- powinna być konfigurowalna.
    \item \textbf{P7.} Parametr \texttt{tlist=[0, 1]} w~\texttt{loss.py} --- brak wyjaśnienia jednostek czasu.
    \item \textbf{P8.} Brak logowania --- tylko \texttt{print()} w~\texttt{generate.py}.
\end{itemize}

% =====================================================================
\section{Kluczowe fragmenty kodu źródłowego}
\label{sec:kod}

\begin{lstlisting}[caption={Klasa bazowa stanów kwantowych (\texttt{state/base.py}).}, label={lst:base-state}]
from abc import ABC, abstractmethod
from qutip import Qobj

class QuantumState(ABC):

    @abstractmethod
    def ket(self) -> Qobj:
        pass

    def density_matrix(self) -> Qobj:
        psi = self.ket()
        return psi * psi.dag()
\end{lstlisting}

\begin{lstlisting}[caption={Klasa bazowa kanałów kwantowych (\texttt{noise/base.py}).}, label={lst:base-channel}]
from abc import ABC, abstractmethod
from qutip import Qobj

class QuantumChannel(ABC):

    @abstractmethod
    def apply(self, rho: Qobj) -> Qobj:
        pass
\end{lstlisting}

\begin{lstlisting}[caption={Implementacja stanu GKP (\texttt{state/gkp.py}).}, label={lst:gkp}]
import numpy as np
from qutip import displace, basis
from .base import QuantumState

class GKPState(QuantumState):

    def __init__(self, cutoff: int, delta: float = 0.3, grid_size: int = 5):
        self.cutoff = cutoff
        self.delta = delta
        self.grid_size = grid_size

    def ket(self):
        psi = 0
        vacuum = basis(self.cutoff, 0)

        for s in range(-self.grid_size, self.grid_size + 1):
            q_shift = 2 * s * np.sqrt(np.pi)
            D = displace(self.cutoff, q_shift / np.sqrt(2))
            weight = np.exp(-(2 * s * np.sqrt(np.pi))**2
                            / (2 * self.delta**2))
            psi += weight * D * vacuum

        return psi.unit()
\end{lstlisting}

\begin{lstlisting}[caption={Generator zbioru danych (\texttt{generate.py}) --- fragment z błędem B2.}, label={lst:generate}]
import os
import numpy as np
import matplotlib.pyplot as plt
from qutip import fock, coherent, wigner, ket2dm

def generate_dataset(n_samples=10, resolution=100):
    paths = setup_folders()
    N = 32
    xvec = np.linspace(-5, 5, resolution)

    for i in range(n_samples):
        # --- FOCK STATE ---
        n_photons = np.random.randint(0, 8)
        state_fock = fock(N, n_photons)
        w_clean = wigner(state_fock, xvec, xvec)
        plt.imsave(os.path.join(paths["clean"],
            f"fock_n{n_photons}_id{i}.png"), w_clean, cmap='RdBu_r')

        # --- BINOMIAL STATE (BLAD! To nie jest stan dwumianowy) ---
        n = np.random.randint(0, 8)
        p = np.random.uniform(0.1, 0.9)
        state_binomial = (fock(N, n) + p * fock(N, n+1)).unit()  # B2!
\end{lstlisting}

% =====================================================================
\section{Plan kolejnych działań --- podsumowanie faz}
\label{sec:plan-podsumowanie}

\begin{table}[htbp]
\centering
\caption{Podsumowanie faz planu działań.}
\label{tab:fazy}
\begin{tabular}{@{}clp{7.5cm}r@{}}
\toprule
\textbf{Faza} & \textbf{Priorytet} & \textbf{Zakres} & \textbf{Czas} \\
\midrule
1 & WYSOKI  & Naprawa błędów krytycznych (B1--B4, weryfikacja wzorów) & $\sim$4h \\
2 & WYSOKI  & Infrastruktura (\texttt{requirements.txt}, parametryzacja, README) & $\sim$5h \\
3 & ŚREDNI  & Jakość kodu (typy, docstringi, walidacja, referencje) & $\sim$11h \\
4 & ŚREDNI  & Testy (pytest, testy stanów, kanałów, pipeline) & $\sim$10h \\
5 & NISKI   & Rozwój (nowe kanały, CI/CD, moduł ML) & otwarte \\
\bottomrule
\end{tabular}
\end{table}

% =====================================================================
\section{Faza 1 --- Naprawa błędów krytycznych}
\label{sec:faza1}

\subsection{Dodanie plików \texttt{\_\_init\_\_.py}}

\textbf{Problem}: Brak plików \texttt{\_\_init\_\_.py} uniemożliwia import pakietów Pythona.

\textbf{Kroki}:
\begin{enumerate}
    \item Utworzyć puste pliki \texttt{\_\_init\_\_.py} w~katalogach: \texttt{src/}, \texttt{src/physics/}, \texttt{src/physics/state/}, \texttt{src/physics/noise/}, \texttt{src/physics/measurement/}.
    \item W~\texttt{src/physics/\_\_init\_\_.py} dodać eksporty głównych klas (Listing~\ref{lst:init}).
    \item Poprawić niespójne importy --- w~\texttt{coherent.py} i~\texttt{vacuum.py} jest import bezwzględny (\texttt{from src.physics.state.base}), a~w~pozostałych plikach względny (\texttt{from .base}). Ujednolicić do importów względnych.
    \item Sprawdzić: \texttt{python -c "from src.physics import FockState"}.
\end{enumerate}

\begin{lstlisting}[caption={Proponowana zawartość \texttt{src/physics/\_\_init\_\_.py}.}, label={lst:init}]
from .state.fock import FockState
from .state.coherent import CoherentState
from .state.cat import CatState
from .state.vacuum import VacuumState
from .state.thermal import ThermalState
from .state.gkp import GKPState
from .state.binomial import BinomialState
from .noise.loss import LossChannel
from .noise.mixture import MixtureChannel
from .measurement.wigner import WignerMeasurement
from .measurement.pipeline import MeasurementPipeline
\end{lstlisting}

\textbf{Kryterium zakończenia}: Polecenie \texttt{python -c "from src.physics import FockState"} wykonuje się bez błędów.

\subsection{Zmiana nazwy \texttt{binomal.py} $\to$ \texttt{binomial.py}}

\begin{enumerate}
    \item \texttt{git mv src/physics/state/binomal.py src/physics/state/binomial.py}
    \item Zaktualizować wszystkie importy.
    \item \texttt{grep -r "binomal"} --- upewnić się, że nie ma starych odwołań.
\end{enumerate}

\subsection{Dodanie dziedziczenia w kanałach szumu}

\begin{lstlisting}[caption={Naprawa dziedziczenia w \texttt{noise/loss.py} i \texttt{noise/mixture.py}.}, label={lst:fix-inheritance}]
# PRZED:
class LossChannel:
# PO:
from .base import QuantumChannel
class LossChannel(QuantumChannel):

# PRZED:
class MixtureChannel:
# PO:
from .base import QuantumChannel
class MixtureChannel(QuantumChannel):
\end{lstlisting}

\textbf{Kryterium}: \texttt{isinstance(LossChannel(32, 0.1), QuantumChannel)} zwraca \texttt{True}.

\subsection{Naprawa generowania stanu dwumianowego}

\begin{lstlisting}[caption={Naprawa błędu B2 w \texttt{generate.py}.}, label={lst:fix-binomial}]
# PRZED (linia 44):
state_binomial = (fock(N, n) + p * fock(N, n+1)).unit()

# PO:
from src.physics.state.binomial import BinomialState
state_binomial = BinomialState(N=n, p=p, cutoff=N).ket()
\end{lstlisting}

\subsection{Weryfikacja wzorów matematycznych}

\subsubsection{Stan dwumianowy}

Obecna implementacja:
\begin{equation}
    \ket{\psi_\mathrm{bin}} = \sum_{n=0}^{N} \sqrt{\binom{N}{n}}\, p^{n/2}\, (1-p)^{(N-n)/2} \ket{n}
    \label{eq:binomial}
\end{equation}

Kroki weryfikacji:
\begin{enumerate}
    \item Sprawdzić definicję w~Michael et~al.~\cite{michael2016}.
    \item Potwierdzić normalizację: $\braket{\psi|\psi} = \sum \binom{N}{n} p^n (1-p)^{N-n} = 1$ (wzór dwumianowy Newtona).
    \item Upewnić się, że \texttt{cutoff > N}.
\end{enumerate}

\subsubsection{Stan GKP}

Idealny stan GKP to nieskończony grzebień delt Diraca~\cite{gottesman2001}:
\begin{align}
    \ket{0_\mathrm{GKP}} &\propto \sum_s \delta(q - 2s\sqrt{\pi}) \label{eq:gkp-ideal-0} \\
    \ket{1_\mathrm{GKP}} &\propto \sum_s \delta(q - (2s+1)\sqrt{\pi}) \label{eq:gkp-ideal-1}
\end{align}

Przybliżenie skończonej energii zastosowane w~projekcie:
\begin{equation}
    \ket{\mathrm{GKP}} \propto \sum_s \exp\!\left(-\frac{(2s\sqrt{\pi})^2}{2\delta^2}\right) \hat{D}\!\left(\frac{2s\sqrt{\pi}}{\sqrt{2}}\right) \ket{0}
    \label{eq:gkp-approx}
\end{equation}

Parametr $\delta$ kontroluje kompromis: mniejszy $\delta$ $\to$ bliżej ideału, ale większa wymagana energia. Na wykresie Wignera stan GKP wyróżnia się charakterystyczną \textbf{regularną siatką pików} (patrz Rysunek~\ref{fig:wigner-examples}).

\textbf{Kroki weryfikacji}: Porównać z~rów.~(18) w~\cite{gottesman2001}. Sprawdzić konwencję QuTiP: $\hat{D}(\alpha)$ przesuwa o~$\mathrm{Re}(\alpha)$ w~$q$.

% =====================================================================
\section{Faza 2 --- Infrastruktura projektu}
\label{sec:faza2}

\subsection{Utworzenie \texttt{requirements.txt}}

\begin{lstlisting}[style=tree, caption={Proponowany \texttt{requirements.txt}.}, label={lst:requirements}]
qutip>=5.0
numpy>=1.24
scipy>=1.10
matplotlib>=3.7
Pillow>=10.0
\end{lstlisting}

Opcjonalnie \texttt{requirements-dev.txt} z~\texttt{pytest>=7.0} i~\texttt{opencv-python>=4.8}.

\subsection{Parametryzacja ścieżek i stałych}

\begin{lstlisting}[caption={Proponowany plik konfiguracyjny \texttt{config.py}.}, label={lst:config}]
from dataclasses import dataclass
from pathlib import Path

@dataclass
class GeneratorConfig:
    output_dir: Path = Path("quantum_dataset")
    cutoff: int = 32
    resolution: int = 100
    x_range: float = 5.0
    n_samples: int = 10
    cmap: str = "RdBu_r"
\end{lstlisting}

Dodać interfejs CLI z~\texttt{argparse}: \texttt{python generate.py --n-samples 100 --resolution 200}.

\subsection{Ujednolicenie języka}

Przetłumaczyć komentarze polskie na angielski (m.in.~w~\texttt{generate.py}, \texttt{gkp.py}).

\textbf{Kryterium}: \texttt{grep -rn "[ąćęłńóśźż]" src/} nie zwraca wyników.

\subsection{Rozbudowa \texttt{README.md}}

Dodać sekcje: opis projektu, struktura, wymagania, instalacja, przykłady użycia, autorzy. Przykład szybkiego startu:

\begin{lstlisting}[caption={Przykład użycia do umieszczenia w README.}, label={lst:quickstart}]
from src.physics.state.fock import FockState
from src.physics.measurement.wigner import WignerMeasurement

state = FockState(n=3, cutoff=32)
wm = WignerMeasurement(x_max=5, resolution=100)
W = wm.measure(state.density_matrix())
\end{lstlisting}

% =====================================================================
\section{Faza 3 --- Jakość kodu}
\label{sec:faza3}

\subsection{Adnotacje typów}

\begin{lstlisting}[caption={Przykład adnotacji typów dla \texttt{FockState}.}, label={lst:types}]
class FockState(QuantumState):
    def __init__(self, n: int, cutoff: int) -> None:
        self.n: int = n
        self.cutoff: int = cutoff

    def ket(self) -> Qobj:
        return basis(self.cutoff, self.n)
\end{lstlisting}

Kryterium: \texttt{mypy src/} przechodzi bez błędów krytycznych (\texttt{mypy} to narzędzie do statycznego sprawdzania poprawności typów w~Pythonie --- wykrywa błędy typów bez uruchamiania kodu).

\subsection{Docstringi}

\begin{lstlisting}[caption={Przykładowy docstring z referencją do literatury (GKP).}, label={lst:docstring-gkp}]
class GKPState(QuantumState):
    """Approximate Gottesman-Kitaev-Preskill (GKP) state.

    GKP state encodes a logical qubit in a harmonic oscillator,
    protecting against small displacement errors in position (q)
    and momentum (p). The ideal state is an infinite comb of
    Dirac deltas spaced by 2*sqrt(pi). This implementation uses
    a finite-energy approximation with a Gaussian envelope of
    width delta:

        |GKP> ~ sum_s exp(-2*pi*s**2 / delta**2) D(s*sqrt(2*pi)) |0>

    Reference:
        D. Gottesman, A. Kitaev, J. Preskill,
        "Encoding a qubit in an oscillator",
        Phys. Rev. A 64, 012310 (2001).
        https://doi.org/10.1103/PhysRevA.64.012310

    Args:
        cutoff: Hilbert space dimension.
        delta: Envelope width (smaller = more ideal).
        grid_size: Number of peaks on each side.
    """
\end{lstlisting}

\subsection{Walidacja danych wejściowych}

\begin{lstlisting}[caption={Przykłady walidacji parametrów.}, label={lst:validation}]
# BinomialState.__init__()
if N < 0:
    raise ValueError(f"N must be non-negative, got {N}")
if not 0 <= p <= 1:
    raise ValueError(f"p must be in [0, 1], got {p}")
if cutoff <= N:
    raise ValueError(f"cutoff ({cutoff}) must be > N ({N})")

# LossChannel.__init__()
if gamma < 0:
    raise ValueError(f"gamma must be non-negative, got {gamma}")
\end{lstlisting}

\subsection{Obsługa błędów w skryptach}

Dodać \texttt{try/except} w~\texttt{generate.py} i~sprawdzanie istnienia pliku w~\texttt{divider.py}.

\subsection{Referencje do literatury}

Dodać referencje w~docstringach: GKP~\cite{gottesman2001}, stan dwumianowy~\cite{michael2016}, kanał strat --- równanie Lindblada~\cite{lindblad1976, wiseman2009}.

% =====================================================================
\section{Faza 4 --- Testy}
\label{sec:faza4}

\subsection{Konfiguracja pytest}

Utworzyć katalog \texttt{tests/} ze strukturą:

\begin{lstlisting}[style=tree, caption={Struktura katalogu testów.}]
tests/
|-- __init__.py
|-- conftest.py          -- wspolne fixtures
|-- test_states.py       -- testy stanow
|-- test_noise.py        -- testy kanalow szumu
|-- test_validation.py   -- testy walidacji
|-- test_measurement.py  -- testy pomiarow
`-- test_pipeline.py     -- testy integracyjne
\end{lstlisting}

\subsection{Testy jednostkowe stanów}

Każdy stan kwantowy musi spełniać warunki fizyczne z~sekcji~\ref{sec:pojecia}. Tabela~\ref{tab:testy-wspolne} przedstawia testy wspólne.

\begin{table}[htbp]
\centering
\caption{Testy wspólne dla wszystkich stanów kwantowych.}
\label{tab:testy-wspolne}
\begin{tabular}{@{}p{3cm}p{3cm}p{4cm}p{4.5cm}@{}}
\toprule
\textbf{Test} & \textbf{Co sprawdzamy?} & \textbf{Dlaczego to ważne?} & \textbf{Kod} \\
\midrule
Normalizacja & $\braket{\psi|\psi} = 1$ & Prawdopodobieństwa = 1 & \texttt{state.ket().norm()} \\
\addlinespace
Ślad $\rho$ & $\mathrm{Tr}(\rho) = 1$ & Jak wyżej, dla macierzy & \texttt{rho.tr()} \\
\addlinespace
Hermitowskość & $\rho = \rho^\dagger$ & Wielkości obserwowalne muszą być rzeczywiste & \texttt{rho.isherm} \\
\addlinespace
Dodatnia półokreśloność & wartości własne $\geq 0$ & Prawdopodobieństwa nieujemne & \texttt{rho.eigenenergies()} \\
\bottomrule
\end{tabular}
\end{table}

Testy specyficzne:

\begin{lstlisting}[caption={Testy stanów kwantowych --- przykłady.}, label={lst:test-states}]
class TestFockState:
    def test_photon_number(self, cutoff):
        """Stan Focka |n> powinien miec dokladnie n fotonow.
        n_hat = a.dag() * a to operator liczby fotonow.
        expect(n_hat, rho) oblicza wartosc oczekiwana <n>.
        """
        state = FockState(n=3, cutoff=cutoff)
        rho = state.density_matrix()
        a = destroy(cutoff)         # operator anihilacji
        n_hat = a.dag() * a         # operator liczby fotonow
        assert abs(expect(n_hat, rho) - 3.0) < 1e-10

class TestCoherentState:
    def test_mean_photon_number(self, cutoff):
        """Stan koherentny |alpha> ma srednio |alpha|^2 fotonow."""
        alpha = 2.0
        state = CoherentState(alpha=alpha, cutoff=cutoff)
        rho = state.density_matrix()
        a = destroy(cutoff)
        n_hat = a.dag() * a
        assert abs(expect(n_hat, rho) - abs(alpha)**2) < 0.01

class TestCatState:
    def test_symmetry(self, cutoff):
        """Parzysty kot zawiera tylko parzyste stany Focka.
        |alpha> + |-alpha> -> tylko |0>, |2>, |4>, ...
        """
        state = CatState(alpha=2.0, cutoff=cutoff)
        psi = state.ket()
        for n in [1, 3, 5]:
            overlap = abs(basis(cutoff, n).dag() * psi)[0][0]
            assert abs(overlap) < 1e-10

class TestThermalState:
    def test_ket_raises(self, cutoff):
        """Stan termiczny jest mieszany -- nie ma ketu."""
        state = ThermalState(n_th=1.0, cutoff=cutoff)
        with pytest.raises(NotImplementedError):
            state.ket()
\end{lstlisting}

\subsection{Testy kanałów szumu}

\begin{infobox}[Zasada testowania kanałów]
Kanał szumu nie może złamać praw fizyki --- wynik musi nadal być poprawną macierzą gęstości (hermitowska, ślad = 1, wartości własne $\geq 0$).
\end{infobox}

\begin{lstlisting}[caption={Testy kanałów szumu --- przykłady.}, label={lst:test-noise}]
class TestLossChannel:
    def test_reduces_photon_number(self, cutoff):
        """Straty zmniejszaja srednia liczbe fotonow."""
        state = FockState(n=5, cutoff=cutoff)
        channel = LossChannel(cutoff=cutoff, gamma=0.3)
        rho_in = state.density_matrix()
        rho_out = channel.apply(rho_in)
        a = destroy(cutoff)
        n_hat = a.dag() * a
        assert expect(n_hat, rho_out) < expect(n_hat, rho_in)

    def test_zero_loss_preserves_state(self, cutoff):
        """Brak szumu (gamma=0) = brak zmian."""
        state = CoherentState(alpha=1.0, cutoff=cutoff)
        channel = LossChannel(cutoff=cutoff, gamma=0.0)
        rho_in = state.density_matrix()
        rho_out = channel.apply(rho_in)
        assert abs((rho_in - rho_out).norm()) < 1e-6
\end{lstlisting}

\subsection{Testy walidacji}

\begin{lstlisting}[caption={Testy funkcji \texttt{is\_physical()} --- wykrywanie macierzy niefizycznych.}, label={lst:test-validation}]
class TestIsPhysical:
    def test_negative_eigenvalue(self, cutoff):
        """Macierz z ujemna wartoscia wlasna jest niefizyczna.
        diag([1.5, -0.5, 0, ...]) -- slad=1, hermitowska,
        ale -0.5 to ujemne prawdopodobienstwo.
        """
        mat = np.diag([1.5, -0.5] + [0.0] * (cutoff - 2))
        rho_bad = Qobj(mat)
        assert not is_physical(rho_bad)
\end{lstlisting}

\subsection{Testy integracyjne pipeline}

\begin{lstlisting}[caption={Test integracyjny: wszystkie stany przez pipeline.}, label={lst:test-pipeline}]
class TestMeasurementPipeline:
    def test_all_states_produce_output(self, cutoff):
        """Kazdy typ stanu daje poprawny obraz Wignera."""
        wm = WignerMeasurement(x_max=5, resolution=32)
        pipeline = MeasurementPipeline(measurement=wm)
        states = [
            FockState(n=1, cutoff=cutoff),
            CoherentState(alpha=1.0, cutoff=cutoff),
            CatState(alpha=2.0, cutoff=cutoff),
            VacuumState(cutoff=cutoff),
            GKPState(cutoff=cutoff, delta=0.4),
            BinomialState(N=3, p=0.5, cutoff=cutoff),
        ]
        for state in states:
            result = pipeline.run(state)
            assert result.shape == (32, 32)
            assert np.isfinite(result).all()
\end{lstlisting}

% =====================================================================
\section{Faza 5 --- Rozwój funkcjonalności}
\label{sec:faza5}

\subsection{Nowe kanały szumu}

Obecnie projekt ma 2~kanały szumu. Proponowane rozszerzenia:

\begin{description}
    \item[Kanał defazowania] --- niszczenie informacji o~fazie (kącie w~przestrzeni fazowej), bez utraty fotonów. Operator kolapsu: $\sqrt{\gamma}\,\hat{n}$. Analogia: światło w~ośrodku o~losowo zmieniającym się współczynniku załamania.
    \item[Kanał wzmocnienia] --- dodawanie fotonów (odwrotność strat). Operator kolapsu: $\sqrt{\gamma}\,\hat{a}^\dagger$. Analogia: wzmacniacz optyczny.
    \item[Kanał depolaryzacji] --- z~prawdopodobieństwem $p$ stan zastępowany ,,białym szumem'': $\rho \to (1-p)\rho + p\cdot \mathbf{I}/d$, gdzie $\mathbf{I}$ to macierz jednostkowa, a~$d$ --- wymiar przestrzeni Hilberta. Analogia: detektor dający losowy wynik.
\end{description}

\subsection{Generowanie stanów mieszanych}

Rozszerzyć generator o~stany z~szumem (Fock + straty, kot + defazowanie). Dodać metadane CSV.

\subsection{Konfiguracja CI/CD}

\begin{lstlisting}[style=yaml, caption={Proponowany \texttt{.github/workflows/tests.yml}.}, label={lst:ci}]
name: Tests
on: [push, pull_request]
jobs:
  test:
    runs-on: ubuntu-latest
    strategy:
      matrix:
        python-version: ["3.10", "3.11", "3.12"]
    steps:
      - uses: actions/checkout@v4
      - uses: actions/setup-python@v5
        with:
          python-version: ${{ matrix.python-version }}
      - run: pip install -r requirements-dev.txt
      - run: pytest tests/ -v --tb=short
\end{lstlisting}

\subsection{Moduł ML --- klasyfikator stanów}

Główny cel projektu. Zadania ML:
\begin{itemize}
    \item \textbf{Klasyfikacja}: obraz Wignera $\to$ typ stanu (Fock / koherentny / kot / GKP / dwumianowy / termiczny),
    \item \textbf{Rekonstrukcja}: zaszumiony obraz $\to$ czysty obraz,
    \item \textbf{Estymacja parametrów}: obraz Wignera $\to$ parametry ($n$, $\alpha$, $p$).
\end{itemize}

\subsection{Optymalizacja wydajności}

Zapamiętywanie (cache) wyników kosztownych obliczeń operatorów przesunięcia w~GKP za pomocą dekoratora \texttt{@lru\_cache} (mechanizm Pythona, który automatycznie zapamiętuje wyniki wywołań funkcji i~zwraca je bez ponownego obliczania). Zamiana pętli na operacje tablicowe NumPy (wektoryzacja --- szybsze niż pętla \texttt{for} w~Pythonie). Równoległe generowanie obrazów za pomocą modułu \texttt{multiprocessing}.

% =====================================================================
\section{Harmonogram i zależności między fazami}
\label{sec:harmonogram}

\begin{figure}[htbp]
\centering
\begin{tikzpicture}[
    phase/.style={draw, rounded corners, minimum width=3.2cm, minimum height=1cm, align=center, font=\small},
    arr/.style={-{Stealth[length=3mm]}, thick},
    node distance=1.5cm and 2cm
]
    \node[phase, fill=red!15] (f1) {Faza 1\\Błędy krytyczne\\$\sim$4\,godz.};
    \node[phase, fill=orange!15, right=of f1] (f2) {Faza 2\\Infrastruktura\\$\sim$5\,godz.};
    \node[phase, fill=yellow!15, right=of f2] (f3) {Faza 3\\Jakość kodu\\$\sim$11\,godz.};
    \node[phase, fill=blue!15, below=of f2] (f4) {Faza 4\\Testy\\$\sim$10\,godz.};
    \node[phase, fill=green!15, right=of f4] (f5) {Faza 5\\Rozwój\\otwarte};

    \draw[arr] (f1) -- (f2);
    \draw[arr] (f2) -- (f3);
    \draw[arr] (f2) -- (f4);
    \draw[arr] (f4) -- (f5);
    \draw[arr, dashed] (f3) -- (f5);
\end{tikzpicture}
\caption{Zależności między fazami planu działań. Linia przerywana = zależność częściowa.}
\label{fig:harmonogram}
\end{figure}

\textbf{Zależności}:
\begin{itemize}
    \item Faza 2 wymaga ukończenia Fazy 1 (pakiety muszą działać przed konfiguracją CLI).
    \item Faza 4 wymaga ukończenia Fazy 1 i~częściowo Fazy 3 (walidacja potrzebna do testów).
    \item Faza 3 i~4 mogą być realizowane częściowo równolegle.
    \item Faza 5 wymaga stabilnej bazy z~Faz 1--4.
\end{itemize}

\textbf{Łączny szacowany nakład}: ok.~30 godzin roboczych (Fazy 1--4) + otwarte (Faza 5, głównie moduł ML).

% =====================================================================
\section{Diagramy architektury}
\label{sec:diagramy}

\begin{figure}[htbp]
\centering
\begin{tikzpicture}[
    class/.style={draw, rectangle, rounded corners=3pt, minimum width=3cm, minimum height=0.8cm, align=center, font=\small\ttfamily},
    abstract/.style={class, fill=blue!10, font=\small\ttfamily\itshape},
    concrete/.style={class, fill=green!10},
    arr/.style={-{Stealth[length=2.5mm]}, thick},
    inherit/.style={-{Triangle[length=3mm, open]}, thick},
    node distance=0.8cm and 1.5cm
]
    % QuantumState hierarchy — all states inherit directly from QuantumState
    \node[abstract] (qs) {QuantumState (ABC)};

    % Row 1: 4 states
    \node[concrete, below left=1.5cm and 3.5cm of qs] (vac) {VacuumState};
    \node[concrete, below left=1.5cm and 1cm of qs] (fock) {FockState};
    \node[concrete, below right=1.5cm and 1cm of qs] (coh) {CoherentState};
    \node[concrete, below right=1.5cm and 3.5cm of qs] (cat) {CatState};
    % Row 2: 3 states
    \node[concrete, below=3cm of qs, xshift=-2.5cm] (therm) {ThermalState};
    \node[concrete, below=3cm of qs] (gkp) {GKPState};
    \node[concrete, below=3cm of qs, xshift=2.5cm] (binom) {BinomialState};

    % All arrows go directly to QuantumState
    \draw[inherit] (vac) -- (qs);
    \draw[inherit] (fock) -- (qs);
    \draw[inherit] (coh) -- (qs);
    \draw[inherit] (cat) -- (qs);
    \draw[inherit] (therm) -- (qs);
    \draw[inherit] (gkp) -- (qs);
    \draw[inherit] (binom) -- (qs);

    % QuantumChannel hierarchy
    \node[abstract, right=4cm of qs] (qc) {QuantumChannel (ABC)};
    \node[concrete, below left=1cm and 0cm of qc] (loss) {LossChannel};
    \node[concrete, below right=1cm and 0cm of qc] (mix) {MixtureChannel};

    \draw[inherit] (loss) -- (qc);
    \draw[inherit] (mix) -- (qc);

    % Note about missing inheritance
    \node[draw=red, dashed, thick, rounded corners, below=0.3cm of loss, font=\scriptsize\color{red}, text width=2.5cm, align=center] {brak dziedziczenia! (błąd B4)};

    % Pipeline
    \node[class, fill=yellow!15, below=3.5cm of qc] (pipe) {MeasurementPipeline};
    \node[class, fill=yellow!15, left=1.5cm of pipe] (wigner) {WignerMeasurement};

    \draw[arr] (pipe) -- node[above, font=\scriptsize]{measurement} (wigner);
    \draw[arr, dashed] (pipe) -- ++(0, 1.5) node[right, font=\scriptsize]{noise (opcjonalnie)} -| ($(qc)+(0,-1)$);
\end{tikzpicture}
\caption{Diagram klas projektu. Wszystkie stany dziedziczą bezpośrednio po \texttt{QuantumState}. Klasy abstrakcyjne (ABC) na niebiesko, konkretne na zielono. Czerwona ramka oznacza wykryty błąd B4.}
\label{fig:architektura}
\end{figure}

\begin{figure}[htbp]
\centering
\begin{tikzpicture}[
    block/.style={draw, rounded corners, minimum width=2.8cm, minimum height=1cm, align=center, font=\small},
    arr/.style={-{Stealth[length=3mm]}, thick},
    node distance=0.5cm and 1.8cm
]
    \node[block, fill=green!15] (state) {Stan kwantowy\\$\ket{\psi}$ / $\rho$};
    \node[block, fill=red!15, right=of state] (noise) {Kanał szumu\\(opcjonalnie)};
    \node[block, fill=blue!15, right=of noise] (measure) {Pomiar\\Wignera};
    \node[block, fill=yellow!15, right=of measure] (image) {Obraz 2D\\(mapa cieplna)};

    \draw[arr] (state) -- node[above, font=\scriptsize]{$\rho$} (noise);
    \draw[arr] (noise) -- node[above, font=\scriptsize]{$\rho'$} (measure);
    \draw[arr] (measure) -- node[above, font=\scriptsize]{$W(q,p)$} (image);
\end{tikzpicture}
\caption{Potok pomiarowy (pipeline): stan kwantowy $\to$ szum $\to$ pomiar Wignera $\to$ obraz.}
\label{fig:pipeline}
\end{figure}

\begin{figure}[htbp]
\centering
\begin{tikzpicture}
    % Placeholder for Wigner function examples
    \foreach \i/\name in {0/Próżnia, 1/Fock $\ket{3}$, 2/Koherentny, 3/Kot, 4/GKP} {
        \begin{scope}[xshift=\i*3.2cm]
            \draw[fill=gray!10, rounded corners] (0,0) rectangle (2.8, 2.8);
            \node[font=\scriptsize, align=center] at (1.4, 1.4) {[obraz Wignera\\do wygenerowania]};
            \node[font=\scriptsize\bfseries, below] at (1.4, 0) {\name};
        \end{scope}
    }
\end{tikzpicture}
\caption{Przykładowe funkcje Wignera dla różnych stanów kwantowych (do wygenerowania z~kodu projektu). Każdy stan ma charakterystyczny, unikalny wzór --- to jest cecha, którą moduł ML będzie klasyfikował.}
\label{fig:wigner-examples}
\end{figure}

% =====================================================================
\section{Podsumowanie}
\label{sec:podsumowanie}

Projekt ma \textbf{solidne fundamenty architektoniczne} --- przejrzysty podział na moduły, poprawne użycie klas abstrakcyjnych i~trafny dobór bibliotek. Większość implementacji fizycznych jest poprawna.

Główne braki to: \textbf{brak struktury pakietu} (\texttt{\_\_init\_\_.py}), \textbf{brak testów}, \textbf{brak dokumentacji} oraz \textbf{kilka błędów matematycznych} do zweryfikowania. Brakuje też części ML, która jest kluczowym celem projektu (klasyfikacja i~rekonstrukcja stanów).

\textbf{Zalecana kolejność prac}: Faza~1 $\to$ Faza~2 $\to$ Fazy~3 i~4 równolegle $\to$ Faza~5 (patrz diagram zależności na Rysunku~\ref{fig:harmonogram}).

\textbf{Szacowany łączny nakład pracy na Fazy 1--4}: ok.~30 godzin roboczych.

% =====================================================================
\printbibliography[title={Bibliografia}]

\end{document}
